\subsection{Axioms and nonrepresentability of the henselized functor}

\begin{proposition}\label{A:hensel-remaining-gate}
\status{Proved}
Let $S=\Spec\F_p[t]$. On affine $S$-schemes put
$P(B)=\alpha_p(B^h_{(t)})$. This functor extends uniquely to an
fppf sheaf on $S$-schemes, commutes with filtered ring colimits,
and satisfies conditions 4--8 of cote 131, scan 14, with actual
local modularity in 7 and locally finite-type tests in 8.
It is separated for arbitrary fpqc covers, but it is not represented
by a separated locally quasi-finite $S$-scheme.
Together with the arbitrary fpqc descent proved below in
Theorem~\ref{A:hensel-fpqc}, these assertions give the complete
counterexample of Theorem~\ref{A:main-counterexample}.
\end{proposition}
\begin{proof}
We first justify the sheaf and colimit assertions. Henselization of
pairs commutes with filtered colimits. Indeed, a finite polynomial
factorization problem and a proposed coprime factorization descend
to a finite stage; a lift is supplied by the henselian property
there. The Jacobson-radical condition also passes to the colimit,
since invertibility of every $1+tb$ is witnessed at a finite stage.
The universal property then identifies the resulting colimit with
the henselization of the colimit pair. Taking the kernel of the
polynomial map $x\mapsto x^p$ commutes with filtered colimits.

For an \'etale faithfully flat map $A\to B$, replace $A$ by
$A^h_{(t)}$. Henselizing $B\otimes_A A^h_{(t)}$ and its tensor
powers gives the same rings as henselizing $B$ and its tensor
powers over $A$: this follows twice from the universal property
of henselization. The equalizer assertion is therefore exactly
de Jong--Olander, Theorem 2, applied to $X=\alpha_p$ over the
henselian base \cite{deJongOlanderHenselian}. This includes finite
Zariski covers. For a finite locally free faithfully flat map,
integral base change gives
\[
 B^h_{(t)}=B\otimes_A A^h_{(t)},\qquad
 (B\otimes_A B)^h_{(t)}
   =(B\otimes_A B)\otimes_A A^h_{(t)}.
\]
Ordinary faithfully flat descent for $\alpha_p$ gives the
equalizer. Every fppf cover is refined by successive Zariski and
finite locally free covers \cite[Lemma 37.48.2]{StacksFppfRefinement}.
The equalizers just proved, stable under base change, imply the
fppf sheaf condition. Zariski gluing supplies the asserted extension
from affines to schemes.

For any faithfully flat $A\to B$, the induced
$A^h_{(t)}\to B^h_{(t)}$ is flat: it is the composite of base
change of the given flat map and an ind-\'etale henselization.
Its closed fibre modulo $t$ is faithfully flat. Since $t$ belongs
to the Jacobson radical of $A^h_{(t)}$, the generization-stable
image of this flat map contains every prime. Thus it is faithfully
flat, which proves injectivity on $P$. This proves the separatedness assertion needed in the later
equalizer argument.

If $A$ is complete noetherian local, then either $t$ is a unit,
in which case $P$ is the terminal functor on $A$-algebras, or
$t$ belongs to its maximal ideal, in which case $(A,tA)$ is
henselian and $P(A)=\alpha_p(A)$. The same alternative applies
to every Artinian quotient, so completeness gives condition 4.
Over an Artinian local base, $t$ is either invertible or nilpotent.
In the latter case every pair $(B,tB)$ is already henselian.
The entire base-changed functor is therefore represented by
$\alpha_{p,A}$, and in the former case by $\Spec A$. This proves
the finite universal Artinian representability in condition 5.
Henselization of a reduced ring is reduced, as follows by taking
its filtered system of \'etale neighborhoods. Hence $P$ has one
element on every valuation ring and every field. Valuative
uniqueness, condition 6, follows.

We now verify the literal actual-modularity premise in 7.
Fix a complete noetherian local $A$. If $t$ is invertible or
nilpotent, the preceding global representations verify 7.
Suppose that $t$ lies in the maximal ideal $\mathfrak m$ but
is not nilpotent. We claim that there is no finite local
$A$-algebra $B$ with an actual locally modular element $\xi$.
There is just one point over every field. The minimal
field-of-definition requirement therefore forces the residue
field of $B$ to be $A/\mathfrak m$.
On all Artinian local $A$-algebras whose closed point lies over
$\mathfrak m$, the functor is exactly $\alpha_p$. Actual
modularity, restricted to these tests, identifies the complete
local representing pairs:
\[
 (B,\xi)\simeq (A[x]/(x^p),x).                 \tag{A.2.2}
\]
Here both algebras are complete because they are finite over $A$.
The identification follows by applying the two universal
properties to their Artinian quotients and passing to the limit;
the two composite maps are identities on every quotient.

Set
\[
 R_0=A[u]_{(\mathfrak m,u)},\qquad
 R=R_0[\epsilon]/(\epsilon^2).
\]
In $R_0^h{}_{(t)}$ there is a unique root
\[
 f^p-f=tu,\qquad f\equiv0\pmod t.
\]
It does not lie in $R_0$. If it did, reduction modulo a minimal
prime $\mathfrak q$ of $A$ avoiding $t$ would give an
Artin--Schreier root of $tu$ in
$\operatorname{Frac}(A/\mathfrak q)(u)$. At infinity $tu$ has
a pole of order one, whereas every pole of $g^p-g$ has order
divisible by $p$. This is impossible. Such a prime exists because
$t$ is not nilpotent.
Integral base change for henselization gives
$R^h_{(t)}=R_0^h{}_{(t)}[\epsilon]/(\epsilon^2)$.
The element $\epsilon f$ belongs to $P(R)$ and has the unique
field-valued reduction. It is not the image of an element of
$\alpha_p(R)$: comparison of the $\epsilon$ coefficients would
put $f$ in $R_0$. The map $R_0\to R_0^h{}_{(t)}$ is faithfully
flat, since $R_0$ is local and $t\in\mathfrak mR_0$, so this
coefficient comparison is valid. This contradicts the universal
property of (A.2.2) on the local test $R$. Our claim follows.
Thus in this third case the premise of 7 never holds, and 7 is
again satisfied. This argument uses actual modularity throughout.

For condition 8, let $S'$ be an irreducible noetherian test base.
If $t$ vanishes at its generic point, it belongs to its
nilradical and is locally nilpotent; then $P_{S'}=\alpha_{p,S'}$.
Otherwise $D(t)$ is a nonempty open containing the generic point,
and $P$ is terminal there. In either case, near the given generic
point the functor has a separated finite-type representation.
For the locally finite-type test $Z$, an isomorphism of local
rings at its modular point spreads to an isomorphism on suitable
open neighborhoods, as proved in Proposition~\ref{A:8-literal}.
This gives the nonempty open required by 8.

Finally suppose $P$ had a separated locally quasi-finite
representation. Base change to $A=\F_p[[t]]$. Its unique closed
fibre point has a finite local open neighborhood over this
henselian base: take a quasi-compact separated quasi-finite
neighborhood, apply Zariski's Main Theorem, and select the local
factor of its finite envelope through that point. An open subset
of a finite local $A$-scheme containing its closed point is the
whole factor. Its coordinate ring with its universal element is
an actual locally modular finite pair, contrary to the third
case of 7 just proved. Hence no such representation exists.
\end{proof}

The fpqc equalizer assertion is proved in Theorem~\ref{A:hensel-fpqc}
by a finite-module argument over strictly henselian local bases.
The failure of strengthened adic effectivity in
Proposition~\ref{A:hensel-adic-failure} is a separate fact; the
counterexample to L2 is Theorem~\ref{A:L2-counterexample}.
